% !PRINTED none

\begin{center}
\gesperrt{Elfter Abschnitt.}
\end{center}

\begin{center}
\gesperrt{Kettenbrüche.}
\end{center}

\begin{center}
\rule{3cm}{0.4pt}
\end{center}

\begin{center}
§. 115.
\end{center}

\begin{center}
Verwandlung rationaler Brüche in Kettenbrüche.
\end{center}

Wenn man eine natürliche Zahl $m$ durch eine andere $n$ nach den gewöhnlichen Regeln dividirt, so erhält man einen Quotienten und einen Rest, der positiv und kleiner als der Divisor $n$ ist. Bedeutet $a$ den Quotienten und $r$ den Rest, so ist
\weq{1}{m = an + r.}
Alle Zahlen $m, m', m''\ldots$, die denselben Rest geben, heissen \gesperrt{restgleich} oder \gesperrt{congruent} nach dem \gesperrt{Modul} $n$, und man drückt dies nach \gesperrt{Gauss} durch das Zeichen aus:
\weq{2}{m \equiv m' \ (\mathrm{mod}\ n).}
Eine solche Formel wird eine \gesperrt{Congruenz} genannt.

Wenn wir also das System der Zahlen
\[
r = 0, 1, 2 \ldots n - 1
\]
betrachten, so ist jede beliebige (positive oder negative) ganze Zahl $m$ \gesperrt{einer} und nur \gesperrt{einer} von diesen Zahlen nach dem Modul $n$ congruent. Dieselbe Eigenschaft hat aber auch jedes System von Zahlen
\[
s = m_0, m_1, m_2 \ldots m_{n-1},
\]
die so ausgewählt sind, dass unter ihren Resten alle Zahlen $r$ (und jede nur einmal) vorkommt. Es ist dann jede beliebige Zahl $m$ einer und nur einer der Zahlen $s$ nach dem Modul $n$
%CONT>>
